% GENERATED FILE. Edit canonical lessons, then run npm run generate:books.
% canonical-source-hash: fnv1a64:3d96d71c45dd3881
% canonical-lessons: ES-C562-volver-a, ES-C562-acabar-por, ES-C562-ponerse-a, ES-C562-echarse-a, ES-C562-alejarse

\chapter{Volver a empezar}
\label{ch:es-volver-a-empezar}
\hlchaptermodality{562}

\begin{chapteropening}
\textbf{By the end of this chapter:} \emph{I can say that someone did something again, ended up doing it, started doing it or burst out doing it, and moved away.}

Complete the last lesson of chapter 562: I can say that someone did something again, ended up doing it, started doing it or burst out doing it, and moved away.
\end{chapteropening}

\section[volver a]{volver a — to do again}

\label{lesson:ES-C562-volver-a}



\begin{quote}
Before the new one: say the Spanish for the return, then the Spanish for the absence.
\end{quote}

\subsection*{You'll want to know: volver a}
\textbf{volver a} — \textquotedblleft{}to do again\textquotedblright{}.

\begin{quote}
Volvió a llamar a las diez.
\end{quote}

\begin{quote}
\emph{She called again at ten.}
\end{quote}

\emph{Volver a} + infinitive: \emph{No volveré a hacerlo} — I won't do it again.

\begin{cousinweb}
\emph{volver a} is built on \emph{volver} 'to return', which is inherited from Latin \emph{volvere} 'to roll, to turn' (Lewis \& Short). To do something again is to turn back to it.
\end{cousinweb}

\begin{grammarlens}[title={What you've built}]
One more word for telling, explaining and describing at B1.
\end{grammarlens}

\subsection*{Your turn}
\begin{itemize}
\raggedright
  \item \emph{Say it:} \emph{volver a}
  \item \emph{Say it:} \emph{volver a}, once more
  \item \emph{Say it:} a sentence of your own with \emph{volver a}
  \item \emph{Recall:} say the Spanish for the return, then the Spanish for the absence, then say \emph{volver a} again
\end{itemize}

\begin{tcolorbox}[breakable,colback=teal!4,colframe=teal!35!black,title={Before you move on}]
What does volver a mean? (\textquotedblleft{}To do again\textquotedblright{}.) Say it once more.
\end{tcolorbox}
\section[acabar por]{acabar por — to end up (doing)}

\label{lesson:ES-C562-acabar-por}



\begin{quote}
Before the new one: say the Spanish for the absence, then the Spanish for to do again.
\end{quote}

\subsection*{You'll want to know: acabar por}
\textbf{acabar por} — \textquotedblleft{}to end up (doing)\textquotedblright{}.

\begin{quote}
Acabó por aceptar el trabajo.
\end{quote}

\begin{quote}
\emph{He ended up accepting the job.}
\end{quote}

\emph{Acabar por} + infinitive, or \emph{acabar} + gerund: \emph{Acabó aceptando.}

\begin{cousinweb}
\emph{acabar por} is built on \emph{acabar} 'to finish', made from \emph{a-} and \emph{cabo} 'an end', which comes from Latin \emph{caput} 'head' (Lewis \& Short). To end up doing something is to reach its end.
\end{cousinweb}

\begin{grammarlens}[title={What you've built}]
One more word for telling, explaining and describing at B1.
\end{grammarlens}

\subsection*{Your turn}
\begin{itemize}
\raggedright
  \item \emph{Say it:} \emph{acabar por}
  \item \emph{Say it:} \emph{acabar por}, once more
  \item \emph{Say it:} a sentence of your own with \emph{acabar por}
  \item \emph{Recall:} say the Spanish for the absence, then the Spanish for to do again, then say \emph{acabar por} again
\end{itemize}

\begin{tcolorbox}[breakable,colback=teal!4,colframe=teal!35!black,title={Before you move on}]
What does acabar por mean? (\textquotedblleft{}To end up (doing)\textquotedblright{}.) Say it once more.
\end{tcolorbox}
\section[ponerse a]{ponerse a — to start (doing), to set about}

\label{lesson:ES-C562-ponerse-a}



\begin{quote}
Before the new one: say the Spanish for to do again, then the Spanish for to end up.
\end{quote}

\subsection*{You'll want to know: ponerse a}
\textbf{ponerse a} — \textquotedblleft{}to start (doing), to set about\textquotedblright{}.

\begin{quote}
Se puso a llover de repente.
\end{quote}

\begin{quote}
\emph{It suddenly started to rain.}
\end{quote}

\emph{Ponerse a} + infinitive: \emph{Me puse a estudiar.}

\begin{cousinweb}
\emph{ponerse a} is built on \emph{poner} 'to put', which is inherited from Latin \emph{pōnere} 'to put, to place' (Lewis \& Short). To start something is to put yourself to it.
\end{cousinweb}

\begin{grammarlens}[title={What you've built}]
One more word for telling, explaining and describing at B1.
\end{grammarlens}

\subsection*{Your turn}
\begin{itemize}
\raggedright
  \item \emph{Say it:} \emph{ponerse a}
  \item \emph{Say it:} \emph{ponerse a}, once more
  \item \emph{Say it:} a sentence of your own with \emph{ponerse a}
  \item \emph{Recall:} say the Spanish for to do again, then the Spanish for to end up, then say \emph{ponerse a} again
\end{itemize}

\begin{tcolorbox}[breakable,colback=teal!4,colframe=teal!35!black,title={Before you move on}]
What does ponerse a mean? (\textquotedblleft{}To start (doing), to set about\textquotedblright{}.) Say it once more.
\end{tcolorbox}
\section[echarse a]{echarse a — to burst out (doing)}

\label{lesson:ES-C562-echarse-a}



\begin{quote}
Before the new one: say the Spanish for to end up, then the Spanish for to start.
\end{quote}

\subsection*{You'll want to know: echarse a}
\textbf{echarse a} — \textquotedblleft{}to burst out (doing)\textquotedblright{}.

\begin{quote}
Al ver la carta, se echó a reír.
\end{quote}

\begin{quote}
\emph{When she saw the letter, she burst out laughing.}
\end{quote}

Mostly with \emph{reír} and \emph{correr}: \emph{Se echó a correr.}

\begin{cousinweb}
\emph{echarse a} is built on \emph{echar} 'to throw', which is inherited from Latin \emph{iactāre} 'to throw about' (Lewis \& Short). To burst out crying is to throw yourself into it.
\end{cousinweb}

\begin{grammarlens}[title={What you've built}]
One more word for telling, explaining and describing at B1.
\end{grammarlens}

\subsection*{Your turn}
\begin{itemize}
\raggedright
  \item \emph{Say it:} \emph{echarse a}
  \item \emph{Say it:} \emph{echarse a}, once more
  \item \emph{Say it:} a sentence of your own with \emph{echarse a}
  \item \emph{Recall:} say the Spanish for to end up, then the Spanish for to start, then say \emph{echarse a} again
\end{itemize}

\begin{tcolorbox}[breakable,colback=teal!4,colframe=teal!35!black,title={Before you move on}]
What does echarse a mean? (\textquotedblleft{}To burst out (doing)\textquotedblright{}.) Say it once more.
\end{tcolorbox}
\section[alejarse]{alejarse — to move away, to go off}

\label{lesson:ES-C562-alejarse}



\begin{quote}
Before the new one: say the Spanish for to start, then the Spanish for to burst out.
\end{quote}

\subsection*{You'll want to know: alejarse}
\textbf{alejarse} — \textquotedblleft{}to move away, to go off\textquotedblright{}.

\begin{quote}
El coche se alejó despacio.
\end{quote}

\begin{quote}
\emph{The car moved slowly away.}
\end{quote}

\emph{Alejarse de} — to move away from.

\begin{cousinweb}
\emph{alejarse} is built on \emph{lejos} 'far', which comes from Latin \emph{laxius} 'more loosely, further apart', from \emph{laxus} 'loose' (Lewis \& Short), with \emph{a-} and \emph{-ar}.
\end{cousinweb}

\begin{grammarlens}[title={What you've built}]
One more word for telling, explaining and describing at B1.
\end{grammarlens}

\subsection*{Your turn}
\begin{itemize}
\raggedright
  \item \emph{Say it:} \emph{alejarse}
  \item \emph{Say it:} \emph{alejarse}, once more
  \item \emph{Say it:} a sentence of your own with \emph{alejarse}
  \item \emph{Recall:} say the Spanish for to start, then the Spanish for to burst out, then say \emph{alejarse} again
\end{itemize}

\begin{tcolorbox}[breakable,colback=teal!4,colframe=teal!35!black,title={Before you move on}]
What does alejarse mean? (\textquotedblleft{}To move away, to go off\textquotedblright{}.) Say it once more.
\end{tcolorbox}
